\subsection{Polynomials with Coefficients in a Noetherian Ring}

Let A be a ring, \(\sigma\) an endomorphism of the ring A, and d an endomorphism of the additive group of A satisfying the relation

\[
d (a b) = \sigma (a) d (b) + d (a) b\tag{1}
\]

for all \(a, b \in A\) . In other words, d is a derivation from the ring A to the (A,A)-bimodule obtained by endowing the additive group of A with the left law of action \((a, x) \mapsto \sigma(a)x\) and the right law of action \((x, a) \mapsto xa\) . We have \(d(1) = 0\) (III, §10, No.~5, p.~557, Proposition 3).

Recall (IV, §1, No.~1, p.~2) that A{[}X{]} denotes the Z-module \(A \otimes_{Z} Z[X]\) of polynomials in one variable with coefficients in A. We endow it with its natural structure of a left A-module. The family \((\mathrm{X}^{n})_{n \in \mathbf{N}}\) is a basis for A{[}X{]} over A. We identify A with its image under the mapping \(a \mapsto a \otimes 1\) .

PROPOSITION 7. --- Let A, \(\sigma\) , \(d\) be as above. There exists a unique ring structure on the group A{[}X{]} with the following properties:

\begin{enumerate}
\def\labelenumi{\alph{enumi})}
\item
  The addition in this ring is the usual addition of A{[}X{]}.
\item
  The multiplication in this ring extends that of A.
\item
  The product in this ring of a sequence \((a,X,\ldots,X)\) , consisting of an element a of A followed by n terms equal to X, is the polynomial \(aX^{n}\) .
\item
  In this ring, we have \(\mathrm{X}a = \sigma(a)\mathrm{X} + d(a)\) for every \(a \in \mathrm{A}\) .
\end{enumerate}

Let E be the endomorphism ring of the additive group A{[}X{]}. The mapping that sends \(a \in A\) to the homothety \(a_{M}\) of the left A-module M = A{[}X{]} is a ring homomorphism from A to E. Consider the elements u, \(\sigma_{M}\) , and \(d_{M}\) of E defined by \(u(\sum b_{n}X^{n}) = \sum b_{n}X^{n+1}\) , \(\sigma_{\mathrm{M}}(\sum b_{n}X^{n}) = \sum \sigma(b_{n})X^{n}\) , \(d_{\mathrm{M}}(\sum b_{n}X^{n}) = \sum d(b_{n})X^{n}\) . For every \(a \in A\) , we have

\[
u a _ {\mathrm{M}} = a _ {\mathrm{M}} u, \quad \sigma_ {\mathrm{M}} a _ {\mathrm{M}} = \sigma (a) _ {\mathrm{M}} \sigma_ {\mathrm{M}}, \quad d _ {\mathrm{M}} a _ {\mathrm{M}} = \sigma (a) _ {\mathrm{M}} d _ {\mathrm{M}} + (d (a) _ {\mathrm{M}}).\tag{2}
\]

\[
X_ {M} = \sigma_ {M} u+ d_ {M} .\tag{3}
\]

It follows from (2) that for every \(a \in \mathbf{A}\) , we have

\[
\mathrm{X} _ {\mathrm{M}} a _ {\mathrm{M}} = \sigma (a) _ {\mathrm{M}} \mathrm{X} _ {\mathrm{M}} + (d (a) _ {\mathrm{M}}).\tag{4}
\]

Consider the mapping \(\varphi : \mathrm{A}[\mathrm{X}] \to \mathrm{E}\) defined by

\[
\varphi \Bigl (\sum a _ {n} \mathrm{X} ^ {n} \Bigr) = \sum (a _ {n}) _ {\mathrm{M}} (\mathrm{X} _ {\mathrm{M}}) ^ {n}.\tag{5}
\]

This is a group homomorphism. An induction argument shows that we have \((\mathrm{X}_{\mathrm{M}})^{n}(1)=\mathrm{X}^{n}\) for every \(n\in N\) . We therefore have \(\varphi(\mathrm{P})(1)=\mathrm{P}\) for every \(P\in\mathrm{A}[X]\) , which proves that the homomorphism \(\varphi\) is injective. We denote its image by B. The set B is a subgroup of E; it contains 1, and it is stable under left multiplication by \(a_{M}\) for \(a\in A\) and by \(X_{M}\) (see (4)). It is therefore a subring of E. The unique ring structure on A{[}X{]} derived from that on B by transfer of structure via \(\varphi\) has the properties of Proposition 7, where property d) results from relation (4).

If A{[}X{]} is endowed with a ring structure that has the properties of Proposition 7, then the left homothety \(\gamma_{X}\) of this ring (I, §8, No.~1, p.~97) necessarily sends \(bX^{n}\) to \(\sigma(b)\mathrm{X}^{n+1} + d(b)\mathrm{X}^{n}\) for \(b \in A\) and \(n \in N\) , hence is equal to \(X_{M}\) . The homothety \(\gamma_{a}\) is necessarily equal to \(a_{M}\) for every \(a \in A\) . Consequently, we have \(\gamma_{\mathrm{P}} = \varphi(\mathrm{P})\) for every \(P \in A[X]\) ; the uniqueness in Proposition 7 follows.

The set A{[}X{]}, endowed with the unique ring structure with the properties of Proposition 7, is denoted by \(A[X]_{\sigma,d}\) and called the polynomial ring in X with coefficients in A, relative to \(\sigma\) and d.~We simply denote it by \(A[X]_{\sigma}\) when d is the zero mapping and by A{[}X{]} when, moreover, \(\sigma\) is the identity mapping on A. This notation is compatible with that introduced in IV, §1, No.~1, p.~1 for a commutative ring A.

Remark. --- The ring \(A[X]_{\sigma,d}\) has the following universal property: given a ring \(A'\) , a ring homomorphism \(f: A \to A'\) , and an element x of \(A'\) such that \(xf(a) = f(\sigma(a))x + f(d(a))\) for every \(a \in A\) , there exists a unique ring homomorphism \(g: A[X]_{\sigma,d} \to A'\) that extends f and maps X onto x.

The uniqueness is clear. Let us therefore show that the mapping \(g: \mathrm{A}[\mathrm{X}]_{\sigma,d} \to \mathrm{A}'\) defined by \(g(\sum a_n\mathrm{X}^n) = \sum f(a_n)x^n\) has the required properties. It extends \(f\) , maps \(\mathrm{X}\) onto \(x\) , and is a group homomorphism. We have \(g(1) = 1\) . For \(a \in \mathrm{A}\) and \(\mathrm{Q} = \sum a_n\mathrm{X}^n\) in \(\mathrm{A}[\mathrm{X}]_{\sigma,d}\) , we have

\[
g (a \mathrm{Q}) = g \left(\sum a a _ {n} X ^ {n}\right) = \sum f (a a _ {n}) x ^ {n} = f (a) \sum f (a _ {n}) x ^ {n} = g (a) g (\mathrm{Q})
\]

as well as

\[
\begin{array}{l} g (\mathrm{XQ}) = g \Bigl (\sum \Bigl (\sigma (a _ {n}) \mathrm{X} ^ {n + 1} + d (a _ {n}) \mathrm{X} ^ {n} \Bigr) \Bigr) \\ = \sum (f (\sigma (a _ {n})) x ^ {n + 1} + f (d (a _ {n})) x ^ {n}) = x \sum f (a _ {n}) x ^ {n} = g (\mathrm{X}) g (\mathrm{Q}). \end{array}
\]

It follows that we have \(g(\mathrm{P})g(\mathrm{Q}) = g(\mathrm{PQ})\) for P, Q in \(A[X]_{\sigma,d}\) and therefore that g is a ring homomorphism.

THEOREM 2. --- Let A be a left Noetherian ring, and let \(\sigma\) be an automorphism of A and \(d\) an endomorphism of the additive group of A satisfying relation (1). The ring \(\mathrm{A}[\mathrm{X}]_{\sigma,d}\) is left Noetherian.

Set \(B = A[X]_{\sigma,d}\) . For any integer \(n \geqslant 0\) , we denote by \(B_{n}\) the set of elements of B of the form \(a_{0} + a_{1}X + \cdots + a_{n}X^{n}\) . It is a left A-submodule of B. The mapping \(\varphi_{n}: B_{n} \to A_{s}\) defined by \(\varphi_{n}(a_{0} + a_{1}X + \cdots + a_{n}X^{n}) = a_{n}\) is A-linear.

Let \(\mathfrak{b}\) be a left ideal of B. For every integer \(n \geqslant 0\) , the set \(\mathfrak{a}_n = \varphi_n(\mathfrak{b} \cap \mathrm{B}_n)\) is a left ideal of A. Since we have \(Xa = \sigma(a)X + d(a)\) for every \(a \in A\) , we have

\[
\varphi_ {n + 1} (\mathrm{XQ}) = \sigma (\varphi_ {n} (\mathrm{Q}))\tag{6}
\]

for every \(Q \in B_{n}\) and therefore \(\sigma(\mathfrak{a}_{n}) \subset \mathfrak{a}_{n+1}\) . Consequently, the sequence \(\mathfrak{a}_{n}^{\prime} = \sigma^{-n}(\mathfrak{a}_{n})\) of ideals of A is increasing. Since the ring A is left Noetherian, there exists an integer \(m \geqslant 0\) such that we have \(\mathfrak{a}_{n}^{\prime} = \mathfrak{a}_{n+1}^{\prime}\) for \(n \geqslant m\) . Since

\(\sigma\) is surjective, we have the relation

\[
\sigma (\mathfrak {a} _ {n}) = \mathfrak {a} _ {n + 1}\tag{7}
\]

for every integer \(n \geqslant m\) .

Let \(\mathfrak{c}\) be the left ideal of B generated by \(\mathfrak{b} \cap \mathrm{B}_m\) . Since the left A-module \(\mathrm{B}_m\) is finitely generated and the ring A is left Noetherian, the left A-module \(\mathfrak{b} \cap \mathrm{B}_m\) is finitely generated (VIII, p.~7, Proposition 4 a)). The left ideal \(\mathfrak{c}\) is therefore generated by a finite subset of B. It is clear that it is contained in \(\mathfrak{b}\) . Let us prove that it is equal to \(\mathfrak{b}\) by proving by induction that for every integer \(n \geqslant 0\) , we have

\[
\mathfrak {b} \cap \mathrm{B} _ {n} \subset \mathfrak {c}.\tag{8}
\]

Relation (8) is true by construction for \(n \leqslant m\) . From now on, we assume that n is an integer \(\geqslant m\) such that \(b \cap B_{n} \subset c\) . Let P be an element of \(b \cap B_{n+1}\) . Then \(\varphi_{n+1}(P)\) belongs to \(\mathfrak{a}_{n+1} = \sigma(\mathfrak{a}_{n})\) , and there consequently exists an element Q of \(b \cap B_{n}\) such that \(\varphi_{n+1}(P) = \sigma(\varphi_{n}(Q))\) . Set R = P - XQ. Because of relation (6), we have \(\varphi_{n+1}(R) = 0\) , that is, \(R \in B_{n}\) . Since P and Q belong to the left ideal b of B, the same is true for R; hence, R and Q belong to \(b \cap B_{n}\) , which is contained in the ideal c by the induction hypothesis. Consequently, P belongs to c.~This proves that we have \(b \cap B_{n+1} \subset c\) .

It follows that \(\mathfrak{b}\) is equal to \(\mathfrak{c}\) ; it is therefore a finitely generated ideal of B. This proves that the ring B is left Noetherian.

If the endomorphism \(\sigma\) of the ring A is not an automorphism, then the ring \(A[X]_{\sigma,d}\) is not necessarily left Noetherian, even when A is a Noetherian commutative ring (VIII, p.~22, Exercise 26).

COROLLARY 1 (Hilbert). --- Let A be a Noetherian commutative ring. For every integer \(n \geqslant 0\) , the polynomial algebra \(A[X_{1}, \ldots, X_{n}]\) is a Noetherian ring.

This follows by induction from Theorem 2, taking Proposition 8 of III, §2, No.~9, p.~453, into account.

COROLLARY 2. --- Let A be a Noetherian commutative ring. A commutative A-algebra generated by finitely many elements is a Noetherian ring.

Such an algebra is isomorphic to an algebra of the form \(A[X_{1},\ldots,X_{n}]/\mathfrak{a}\) , where \(n \geqslant 0\) and a is an ideal of \(A[X_{1},\ldots,X_{n}]\) . We then apply Corollary 1 and Proposition 5 of VIII, p.~7.

COROLLARY 3. --- Every commutative ring is the union of a right directed family of Noetherian subrings.

Indeed, let A be a commutative ring. The subrings of A generated (as Z-algebras) by finitely many elements are Noetherian by Corollary 2. They form a right directed family of subrings of A, with union A.

